%%%%%%%%%%%%%%%%
%%% exod 5 %%%
%%%%%%%%%%%%%%%%

\Chapter{5}

\Verse{1}
{
	\hRJE{וְאַחַר בָּאוּ מֹשֶׁה וְאַהֲרֹן וַיֹּאמְרוּ אֶל פַּרְעֹה כֹּה אָמַר יְהֹוָה אֱלֹהֵי יִשְׂרָאֵל שַׁלַּח אֶת עַמִּי וְיָחֹגּוּ לִי בַּמִּדְבָּר׃}
}
{
	\eRJE{
		5 And after that Moses and Aaron came and said to Pharaoh,
		“YHWH, {\God} of Israel, said this: ‘Let my people go, so
		they will celebrate a festival for me in the wilderness.’”
	}
}
{
%	\footnote{
%		Moses and Aaron came and said to Pharaoh. It is against the background
%		of the confrontation between Moses and God at the bush that we must
%		picture the confrontation between Moses and the Pharaoh. The narrative
%		that precedes the burning-bush scene gives us relatively little
%		background of the sort that would add depth and understanding of the
%		character of Moses at the point of his meeting God. Moses’ meeting with
%		the Pharaoh, on the other hand, comes after we have seen him in
%		conversation with the deity. That scene has revealed facets of his
%		character, both strengths and weaknesses, that are now a part of the
%		figure of Moses as he arrives on the scene of human affairs. Moreover,
%		the encounter with God can hardly be disregarded itself as having
%		impacted upon Moses’ personality. He has received more than the three
%		miracles with which to impress the Israelites and the Pharaoh. He has
%		spoken with God. Presumably, setting foot on the floor of the palace is
%		different after one has set foot on holy ground, and conversing with a
%		king is different after one has conversed with the creator. Exodus leads
%		us through a series of dynamics as Moses confronts, first, the deity,
%		then a king, and then a community, with each dynamic growing more
%		complicated as it becomes richer in background.
%	}
}

\Verse{2}
{
	\hRJE{וַיֹּאמֶר פַּרְעֹה מִי יְהֹוָה אֲשֶׁר אֶשְׁמַע בְּקֹלוֹ לְשַׁלַּח אֶת יִשְׂרָאֵל לֹא יָדַעְתִּי אֶת יְהֹוָה וְגַם אֶת יִשְׂרָאֵל לֹא אֲשַׁלֵּחַ׃}
}
{
	\eRJE{
		And Pharaoh said, “Who is YHWH that I should listen to His
		voice, to let Israel go?! I don’t know YHWH, and also I
		won’t let Israel go.”
	}
}
{
%	\footnote{
%		Who is YHWH. The issue of YHWH’s becoming known, which is to become a
%		major component of subsequent books, begins in Exodus. In Genesis, YHWH
%		is known personally only to a few individuals and to no nations. God’s
%		actions in the world are not identifiably God’s. That is, the narrator
%		informs the reader that it is YHWH who brings the flood, but the masses
%		of humans in the narrative are not portrayed as being aware of the
%		source of their catastrophe. There is likewise no suggestion that the
%		generation of the tower of Babylon know whose force is confounding and
%		scattering them. The same is true of the victims of the destruction of
%		Sodom and Gomorrah, and the only eyewitness turns to salt. Joseph
%		informs the Pharaoh that it is God who brings his dream, who makes its
%		interpretation known, and who is the cause of the events that it
%		reveals; but Joseph refers to the deity only generically, as “God.” He
%		is never pictured as revealing to the Egyptians the name of the deity
%		who is at work (Gen 41:16,25,28,32). In Exodus, however, YHWH makes His
%		presence known to Israel, to Egypt, and to other inhabitants of the
%		region (Exod 15:14–15). The Egyptian king’s first words to Moses are,
%		“Who is YHWH . . . ? I don’t know YHWH” (5:2). By the end of the story
%		he knows.
%	}
}

\Verse{3}
{
	\hRJE{וַיֹּאמְרוּ אֱלֹהֵי הָעִבְרִים נִקְרָא עָלֵינוּ נֵלְכָה נָּא דֶּרֶךְ שְׁלֹשֶׁת יָמִים בַּמִּדְבָּר וְנִזְבְּחָה לַיהֹוָה אֱלֹהֵינוּ פֶּן יִפְגָּעֵנוּ בַּדֶּבֶר אוֹ בֶחָרֶב׃}
}
{
	\eRJE{
		And they said, “The {\God} of the Hebrews has communicated
		with us. Let us go on a trip of three days in the
		wilderness so we may sacrifice to YHWH, our {\God}, or else
		He’ll strike us with an epidemic or with the sword.”
	}
}
{
%	\footnote{
%		Let us go on a trip of three days. Moses and Aaron do not ask for the
%		Israelites’ liberation from Egypt. They ask only to make a three-day
%		sacrificial pilgrimage. They just do not mention the fact that they do
%		not plan to return! Footnote 5:3. strike. The word here (Hebrew pg‘) is
%		the same as the word for the Israelites’ “meeting” Moses later in v. 20.
%		The range of meanings of this root makes a wordplay that is enhanced by
%		clustering with the blatant pun in the next verse. Footnote 5:3. strike
%		us. Moses and Aaron say that God may strike “us,” the Israelites, if
%		they don’t go to sacrifice to Him, but the Hebrew suffix has a subtle
%		double meaning: it can also mean “strike him”—which could apply to
%		Pharaoh himself or to a nonspecific person, that is, to the Egyptians
%		who will in fact suffer the plagues. The double meaning here is not
%		merely playful but meaningful to the context.
%	}
}

\Verse{4}
{
	\hRJE{וַיֹּאמֶר אֲלֵהֶם מֶלֶךְ מִצְרַיִם לָמָּה מֹשֶׁה וְאַהֲרֹן תַּפְרִיעוּ אֶת הָעָם מִמַּעֲשָׂיו לְכוּ לְסִבְלֹתֵיכֶם׃}
}
{
	\eRJE{
		And the king of Egypt said to them, “Why, Moses and Aaron,
		do you turn the people loose from its work? Go to your
		burdens!”
	}
}
{
%	\footnote{
%		turn loose. With the Hebrew root pr‘, this is a pun on “Pharaoh.” An
%		even stronger play on this word will occur later, in the middle of the
%		golden calf episode (Exod 32:25). Footnote 5:4. Go to your burdens! He
%		talks down to Moses and Aaron, telling them to get back to work, and
%		thus conveying to them that they are just slaves like everyone else.
%	}
}

\Verse{5}
{
	\hRJE{וַיֹּאמֶר פַּרְעֹה הֵן רַבִּים עַתָּה עַם הָאָרֶץ וְהִשְׁבַּתֶּם אֹתָם מִסִּבְלֹתָם׃}
}
{
	\eRJE{
		And Pharaoh said, “Here, the people of the land are now
		many, and you’ve made them cease from their burdens.”
	}
}
{
%	\footnote{
%		cease. This is the term that describes God’s ceasing on the seventh day
%		of creation, Hebrew bt, cognate to the word “Sabbath.”
%	}
}

\Verse{6}
{
	\hRJE{וַיְצַו פַּרְעֹה בַּיּוֹם הַהוּא אֶת הַנֹּגְשִׂים בָּעָם וְאֶת שֹׁטְרָיו לֵאמֹר׃}
}
{
	\eRJE{
		And Pharaoh commanded the taskmasters among the people and
		its officers in that day, saying,
	}
}
{
}

\Verse{7}
{
	\hRJE{לֹא תֹאסִפוּן לָתֵת תֶּבֶן לָעָם לִלְבֹּן הַלְּבֵנִים כִּתְמוֹל שִׁלְשֹׁם הֵם יֵלְכוּ וְקֹשְׁשׁוּ לָהֶם תֶּבֶן׃}
}
{
	\eRJE{
		“You shall not continue to give straw to the people to
		make bricks as yesterday and the day before. They shall go
		and collect straw for themselves.
	}
}
{
}

\Verse{8}
{
	\hRJE{וְאֶת מַתְכֹּנֶת הַלְּבֵנִים אֲשֶׁר הֵם עֹשִׂים תְּמוֹל שִׁלְשֹׁם תָּשִׂימוּ עֲלֵיהֶם לֹא תִגְרְעוּ מִמֶּנּוּ כִּי נִרְפִּים הֵם עַל כֵּן הֵם צֹעֲקִים לֵאמֹר נֵלְכָה נִזְבְּחָה לֵאלֹהֵינוּ׃}
}
{
	\eRJE{
		And you shall impose on them the quota of the bricks that
		they were making yesterday and the day before. You shall
		not subtract from it. Because they’re lazy. On account of
		this they’re crying, saying, ‘Let us go and sacrifice to
		our {\God}.’
	}
}
{
}

\Verse{9}
{
	\hRJE{תִּכְבַּד הָעֲבֹדָה עַל הָאֲנָשִׁים וְיַעֲשׂוּ בָהּ וְאַל יִשְׁעוּ בְּדִבְרֵי שָׁקֶר׃}
}
{
	\eRJE{
		Let the work be heavy on the people, and let them do it,
		and let them not pay attention to words that are a lie!”
	}
}
{
%	\footnote{
%		Let the work be heavy. The idiom of Moses’ “heavy mouth and heavy
%		tongue” initiates a chain of puns on the various shades of meaning of
%		the word “heavy” (Hebrew kbd, meaning weighty, difficult, or
%		substantial). Here Pharaoh orders that the work of the Israelites be
%		made heavy. The hardening of the Pharaoh’s heart is expressed several
%		times through this term (7:14; 8:11,28; 9:7,34; 10:1). Four of the
%		plagues on Egypt are described as heavy: insects (8:20), pestilence
%		(9:3), hail (9:18,24), and locusts (10:14). The Israelites leave with a
%		heavy supply of livestock (12:38). As Moses holds up his arms in support
%		of the Israelites while they battle the Amalekites, his arms become
%		heavy (17:12). Moses’ father-in-law warns him that his efforts to manage
%		the entire nation on his own are foolish, “because the thing is too
%		heavy for you” (18:18). The chain of puns culminates at Sinai, as a
%		heavy cloud is on the mountain while the people hear the voice of God
%		(19:16). It is not that the narrator’s vocabulary is too limited to
%		produce other adjectives. Rather, the construction of such chains of
%		wordplay is a technique of biblical narrative, deployed in other
%		biblical books as well. Here the chain contributes to the narrative’s
%		unity, rather than the narrative’s appearing to be a series of loosely
%		bound episodes.
%	}
}

\Verse{10}
{
	\hRJE{וַיֵּצְאוּ נֹגְשֵׂי הָעָם וְשֹׁטְרָיו וַיֹּאמְרוּ אֶל הָעָם לֵאמֹר כֹּה אָמַר פַּרְעֹה אֵינֶנִּי נֹתֵן לָכֶם תֶּבֶן׃}
}
{
	\eRJE{
		And the people’s taskmasters and its officers went out and
		said to the people, saying, “Pharaoh said this: ‘I am not
		giving you straw.
	}
}
{
}

\Verse{11}
{
	\hRJE{אַתֶּם לְכוּ קְחוּ לָכֶם תֶּבֶן מֵאֲשֶׁר תִּמְצָאוּ כִּי אֵין נִגְרָע מֵעֲבֹדַתְכֶם דָּבָר׃}
}
{
	\eRJE{
		You, go, take straw for yourselves from wherever you’ll
		find it, because not a thing is subtracted from your
		work.’”
	}
}
{
}

\Verse{12}
{
	\hRJE{וַיָּפֶץ הָעָם בְּכל אֶרֶץ מִצְרָיִם לְקֹשֵׁשׁ קַשׁ לַתֶּבֶן׃}
}
{
	\eRJE{
		And the people scattered through all the land of Egypt to
		collect stubble for straw.
	}
}
{
}

\Verse{13}
{
	\hRJE{וְהַנֹּגְשִׂים אָצִים לֵאמֹר כַּלּוּ מַעֲשֵׂיכֶם דְּבַר יוֹם בְּיוֹמוֹ כַּאֲשֶׁר בִּהְיוֹת הַתֶּבֶן׃}
}
{
	\eRJE{
		And the taskmasters were prodding, saying, “Finish your
		work, the day’s thing in that very day, as when there was
		the straw.”
	}
}
{
}

\Verse{14}
{
	\hRJE{וַיֻּכּוּ שֹׁטְרֵי בְּנֵי יִשְׂרָאֵל אֲשֶׁר שָׂמוּ עֲלֵהֶם נֹגְשֵׂי פַרְעֹה לֵאמֹר מַדּוּעַ לֹא כִלִּיתֶם חקְכֶם לִלְבֹּן כִּתְמוֹל שִׁלְשֹׁם גַּם תְּמוֹל גַּם הַיּוֹם׃}
}
{
	\eRJE{
		And the officers of the children of Israel whom Pharaoh’s
		taskmasters had set over them were beaten—saying, “Why
		haven’t you finished your requirement to make brick like
		the day before yesterday, also yesterday, also today?”
	}
}
{
}

\Verse{15}
{
	\hRJE{וַיָּבֹאוּ שֹׁטְרֵי בְּנֵי יִשְׂרָאֵל וַיִּצְעֲקוּ אֶל פַּרְעֹה לֵאמֹר לָמָּה תַעֲשֶׂה כֹה לַעֲבָדֶיךָ׃}
}
{
	\eRJE{
		And the officers of the children of Israel came and cried
		to Pharaoh, saying, “Why do you do a thing like this to
		your servants?
	}
}
{
}

\Verse{16}
{
	\hRJE{תֶּבֶן אֵין נִתָּן לַעֲבָדֶיךָ וּלְבֵנִים אֹמְרִים לָנוּ עֲשׂוּ וְהִנֵּה עֲבָדֶיךָ מֻכִּים וְחָטָאת עַמֶּךָ׃}
}
{
	\eRJE{
		Straw isn’t given to your servants, and they say to us:
		Make bricks. And here, your servants are beaten, and it’s
		your people’s sin.”
	}
}
{
}

\Verse{17}
{
	\hRJE{וַיֹּאמֶר נִרְפִּים אַתֶּם נִרְפִּים עַל כֵּן אַתֶּם אֹמְרִים נֵלְכָה נִזְבְּחָה לַיהֹוָה׃}
}
{
	\eRJE{
		And he said, “You’re lazy, lazy! On account of this you’re
		saying, ‘Let us go and sacrifice to YHWH.’
	}
}
{
}

\Verse{18}
{
	\hRJE{וְעַתָּה לְכוּ עִבְדוּ וְתֶבֶן לֹא יִנָּתֵן לָכֶם וְתֹכֶן לְבֵנִים תִּתֵּנוּ׃}
}
{
	\eRJE{
		And now go, work! And straw will not be given to you, and
		you shall give the quota of bricks!”
	}
}
{
}

\Verse{19}
{
	\hRJE{וַיִּרְאוּ שֹׁטְרֵי בְנֵי יִשְׂרָאֵל אֹתָם בְּרָע לֵאמֹר לֹא תִגְרְעוּ מִלִּבְנֵיכֶם דְּבַר יוֹם בְּיוֹמוֹ׃}
}
{
	\eRJE{
		And the officers of the children of Israel saw themselves
		in a bad state—saying, “You shall not subtract from your
		bricks, the day’s thing in that very day.”
	}
}
{
}

\Verse{20}
{
	\hRJE{וַיִּפְגְּעוּ אֶת מֹשֶׁה וְאֶת אַהֲרֹן נִצָּבִים לִקְרָאתָם בְּצֵאתָם מֵאֵת פַּרְעֹה׃}
}
{
	\eRJE{
		And they met Moses and Aaron, standing opposite them when
		they came out from Pharaoh,
	}
}
{
}

\Verse{21}
{
	\hRJE{וַיֹּאמְרוּ אֲלֵהֶם יֵרֶא יְהֹוָה עֲלֵיכֶם וְיִשְׁפֹּט אֲשֶׁר הִבְאַשְׁתֶּם אֶת רֵיחֵנוּ בְּעֵינֵי פַרְעֹה וּבְעֵינֵי עֲבָדָיו לָתֶת חֶרֶב בְּיָדָם לְהרְגֵנוּ׃}
}
{
	\eRJE{
		and they said to them, “May YHWH look on you and judge,
		that you’ve made our smell odious in Pharaoh’s eyes and in
		his servants’ eyes, to give a sword in their hand to kill
		us!”
	}
}
{
%	\footnote{
%		our smell odious in Pharaoh’s eyes. People smell with their noses, not
%		their eyes! This may be just an expression or even the author’s
%		oversight, but I rather think that the contradiction is purposeful: the
%		people’s officers misspeak because they are upset. We frequently do
%		that.
%	}
}

\Verse{22}
{
	\hRJE{וַיָּשׁב מֹשֶׁה אֶל יְהֹוָה וַיֹּאמַר אֲדֹנָי לָמָה הֲרֵעֹתָה לָעָם הַזֶּה לָמָּה זֶּה שְׁלַחְתָּנִי׃}
}
{
	\eRJE{
		And Moses went back to YHWH and said, “My Lord, why have
		you done bad to this people? Why did you send me here?
	}
}
{
}

\Verse{23}
{
	\hRJE{וּמֵאָז בָּאתִי אֶל פַּרְעֹה לְדַבֵּר בִּשְׁמֶךָ הֵרַע לָעָם הַזֶּה וְהַצֵּל לֹא הִצַּלְתָּ אֶת עַמֶּךָ׃}
}
{
	\eRJE{
		And since I came to Pharaoh to speak in your name he has
		done bad to this people, and you haven’t rescued your
		people!”
	}
}
{
}
